\section*{Hoofdstuk \ref{ch:g12:buffers-predominance} --- Buffers en predominantiediagrammen}

\begin{solution}{exo:g12:buffers-predominance:1}
\omterm{def:g9:acids-bases-ph:ph}{pH} 2: \ce{CH3COOH}. \omterm{def:g9:acids-bases-ph:ph}{pH} 4.76: beide in gelijke hoeveelheden. \omterm{def:g9:acids-bases-ph:ph}{pH} 7:
\ce{CH3COO-}.
\end{solution}

\begin{solution}{exo:g12:buffers-predominance:2}
\omterm{def:g9:acids-bases-ph:ph}{pH} 4: $1/(1 + 10^{0.76}) = 0.15$. \omterm{def:g9:acids-bases-ph:ph}{pH} 6: $1/(1 + 10^{-1.24}) = 0.95$.
\end{solution}

\begin{solution}{exo:g12:buffers-predominance:3}
Geel; groen (binnen het \omterm{def:g12:buffers-predominance:indicator}{omslagtraject}); blauw.
\end{solution}

\begin{solution}{exo:g12:buffers-predominance:4}
Bij \omterm{def:g9:acids-bases-ph:ph}{pH} 7.4, onder 9.26: \ce{NH4+}. Bij \omterm{def:g9:acids-bases-ph:ph}{pH} 11: \ce{NH3}.
\end{solution}

\begin{solution}{exo:g12:buffers-predominance:5}
Een indicator is zelf een \omterm{def:g12:ka-and-pka:bronsted}{zuur-basekoppel}: in grote hoeveelheden zou hij
de \omterm{def:g9:acids-bases-ph:ph}{pH} veranderen die hij moet laten zien.
\end{solution}

\begin{solution}{exo:g12:buffers-predominance:6}
$10^{\mathrm{pH} - 4.76}$: 0.1; 1; 10.
\end{solution}

\begin{solution}{exo:g12:buffers-predominance:7}
Fenolftaleïne (8.0--10.0) of thymolblauw (8.0--9.1): hun \omterm{def:g12:buffers-predominance:indicator}{omslagtrajecten}
bevatten 8.7. Methylrood slaat om tussen 4.2 en 6.3, ver van 8.7: het zou
al lang eerder omgeslagen zijn.
\end{solution}

\begin{solution}{exo:g12:buffers-predominance:8}
$\mathrm{pH} = 4.76 + \log(0.10/0.20) = 4.76 - 0.30 = 4.46$.
\end{solution}

\begin{solution}{exo:g12:buffers-predominance:9}
\omterm{def:g9:acids-bases-ph:ph}{pH} 1: het \omterm{def:g9:ions:ion}{kation} \ce{H3N+-CH2-COOH}, lading $+1$. \omterm{def:g9:acids-bases-ph:ph}{pH} 6: het zwitterion,
totale lading 0. \omterm{def:g9:acids-bases-ph:ph}{pH} 12: het \omterm{def:g9:ions:ion}{anion} \ce{H2N-CH2-COO-}, lading $-1$.
\end{solution}

\begin{solution}{exo:g12:buffers-predominance:10}
Water: van 7.0 naar $-\log(\num{1.0e-4}/0.101) = 3.0$, een daling van 4
eenheden. Buffer: van 4.76 naar $4.76 + \log(9.9/10.1) = 4.75$, een
daling van 0.01.
\end{solution}

\begin{solution}{exo:g12:buffers-predominance:11}
\ce{NH4+}/\ce{NH3} ($pK_a = 9.26$). Verhouding
$[\ce{NH3}]/[\ce{NH4+}] = 10^{9.5 - 9.26} = 1.7$.
\end{solution}

\begin{solution}{exo:g12:buffers-predominance:12}
De \omterm{def:g9:acids-bases-ph:ph}{pH} daalt met één eenheid als de verhouding $1/10$ bereikt:
$(10 - n)/(10 + n) = 0.1$, dus $n = \qty{8.2}{mmol}$. Een buffer met meer
van elke vorm kan meer zuur of base opvangen voor dezelfde verandering van
de verhouding: hij heeft een grotere capaciteit.
\end{solution}

\begin{solution}{exo:g12:buffers-predominance:13}
Verhouding $10^{5.0 - 4.76} = 1.74$; het zuur is \qty{10}{mmol}, dus
\qty{17.4}{mmol} ethanoaat, dat is \qty{174}{mL} van de oplossing.
\end{solution}

\begin{solution}{exo:g12:buffers-predominance:14}
Het heeft twee \omterm{def:g12:ka-and-pka:bronsted}{zuur-basekoppels}, met twee $pK_a$-waarden, en dus twee
kleuromslagen: rood bij \omterm{def:g9:acids-bases-ph:ph}{pH} 1, geel bij \omterm{def:g9:acids-bases-ph:ph}{pH} 5, blauw bij \omterm{def:g9:acids-bases-ph:ph}{pH} 10.
\end{solution}

\begin{solution}{exo:g12:buffers-predominance:15}
Verdunnen deelt beide concentraties door 10: de verhouding, en dus de \omterm{def:g9:acids-bases-ph:ph}{pH}
(4.76), verandert niet. Een \omterm{def:g12:ka-and-pka:strong-acid}{sterk zuur} dat tien keer verdund wordt, stijgt
één pH-eenheid.
\end{solution}

\begin{solution}{pb:g12:buffers-predominance:1}
\textbf{1.} \ce{CO2 + 2H2O <=> HCO3- + H3O+}.

\textbf{2.} Opgelost koolstofdioxide (met water) is het zuur,
waterstofcarbonaat de base.

\textbf{3.} De tabelwaarde geldt bij \qty{25}{\celsius} in zuiver water;
bloed is \qty{37}{\celsius} en zit vol andere \omterm{def:g9:ions:ion}{ionen}, die de constante
veranderen (en fysiologen tellen al het opgeloste koolstofdioxide mee).

\textbf{4.} Bij \qty{24}{mmol/L}: $0.024 \times 5.0 = \qty{0.12}{mol}$.

\textbf{5.} Een pH-as die bij 6.1 is gedeeld: \ce{CO2} eronder,
\ce{HCO3-} erboven.

\textbf{6.} Bij 7.4, boven 6.1: waterstofcarbonaat.

\textbf{7.} Licht basisch: de \omterm{def:g9:acids-bases-ph:ph}{pH} ligt boven 7.

\textbf{8.} $10^{7.4 - 6.1} = 10^{1.3} = 20$.

\textbf{9.} $24 / 20 = \qty{1.2}{mmol/L}$.

\textbf{10.} $10^{1.28} = 19$ en $10^{1.32} = 21$.

\textbf{11.} $20/21 = \qty{95}{\percent}$.

\textbf{12.} $10^{7.6 - 6.1} = 10^{1.5} = 32$.

\textbf{13.} Het waterstofcarbonaation:
\ce{HCO3- + H3O+ -> CO2 + 2H2O}.

\textbf{14.} $10^{7.1 - 6.1} = 10$.

\textbf{15.} Van $24/20 = 1.2$ naar $24/10 = \qty{2.4}{mmol/L}$:
verdubbeld. De longen ademen sneller en dieper, om het extra
koolstofdioxide uit te blazen.

\textbf{16.} Een \omterm{def:g12:ka-and-pka:strong-acid}{zwak zuur} en zijn base kunnen elk opvangen wat er wordt
toegevoegd, en de \omterm{def:g9:acids-bases-ph:ph}{pH} hangt alleen af van hun verhouding; een \omterm{def:g12:ka-and-pka:strong-acid}{sterk zuur}
zou de \omterm{def:g9:acids-bases-ph:ph}{pH} alleen maar één kant op duwen.

\textbf{17.} Helemaal niet: de verhouding, en dus de \omterm{def:g9:acids-bases-ph:ph}{pH}, blijft gelijk.

\textbf{18.} Ongeveer 20.
\end{solution}
